\documentclass[a4paper,11pt]{article}
\usepackage[utf8]{inputenc}
\usepackage[T1]{fontenc}
\usepackage[french]{babel}
\usepackage{amsmath,amsfonts,amssymb}
\usepackage{geometry}
\usepackage{fancyhdr}
\usepackage{titlesec}
\usepackage{enumitem}
\usepackage{listings}

\geometry{margin=2.5cm}

% Header and footer configuration
\pagestyle{fancy}
\fancyhf{}
\fancyhead[L]{Licence 1 Informatique}
\fancyhead[R]{CERI - Avignon Université}
\fancyfoot[C]{Mathématiques Discrètes \hfill TD1 - Manipuler la logique \hfill Page \thepage/3}

% Title formatting
\titleformat{\section}
{\Large\bfseries}
{\thesection}
{1em}
{}

\titleformat{\subsection}
{\large\bfseries}
{\thesubsection}
{1em}
{}

% Configuration listings pour pseudo-code
\lstset{
    basicstyle=\ttfamily\small,
    keywordstyle=\bfseries,
    commentstyle=\itshape\color{gray},
    frame=single,
    breaklines=true,
    showstringspaces=false
}

\begin{document}

\begin{center}
{\Huge\bfseries TD1 - Manipuler la logique} \\[0.5cm]
{\large Hiver 2026} \\[0.3cm]
\end{center}

\vspace{0.5cm}

L'objectif de ce TD est de manipuler les concepts de base de la logique afin de les prendre en main pour les maîtriser.

\section{Formules et priorité}

Ajoutez des parenthèses aux formules à priorités ci-dessous afin de les transformer en formules strictes équivalentes.

\begin{enumerate}
\item $p \vee q \vee r \vee s \vee t$
\item $p \wedge q \wedge r \wedge s \wedge t$
\item $p \Leftrightarrow \neg q \vee r$
\item $p \vee q \Rightarrow r \wedge s$
\item $p \vee q \Rightarrow r \Leftrightarrow s$
\item $p \vee q \wedge r \Rightarrow \neg s$
\item $p \Rightarrow r \wedge s \Rightarrow t$
\item $p \vee q \wedge s \vee t$
\item $p \wedge q \Leftrightarrow \neg r \vee s$
\item $\neg p \wedge q \vee r \Rightarrow s \Leftrightarrow t$
\end{enumerate}

\section{Logique et français}

Traduisez les phrases en français ci-dessous en formules logiques en précisant quelles variables sont utilisées pour quelles propositions.

\begin{enumerate}
\item Il est faux de dire que nous sommes des géants.
\item Nous sommes inscrits à l'université mais pas à AMU.
\item Les étudiants qui ont suivi des cours de mathématiques ou d'informatique peuvent s'inscrire.
\item Je vais au ski cet hiver ou je voyage à l'étranger cet été, mais pas les deux.
\item Il est nécessaire et suffisant d'être majeur pour avoir le droit de voter.
\item Je suis riche seulement si je gagne à la loterie
\item Si le temps est clément et que je suis en forme, alors je vais me balader
\item Aujourd'hui n'est pas un jour férié.
\item Je suis des cours de mathématiques et des cours d'informatique.
\item S'il pleut, j'irai faire les courses. S'il ne pleut pas, j'irai faire les courses quand même, et j'irai en plus me balader à la plage.
\end{enumerate}

\section{Implication logique et français}

Pour chacune des phrases en français suivantes, donnez la phrase correspondant à la proposition inverse, la proposition réciproque et la proposition contraposée. Précisez quelles phrases ont un sens équivalent.

\begin{enumerate}
\item S'il pleut, je vais faire les courses.
\item J'ai mon diplôme seulement si je valide mon année.
\item Si le temps est clément et que je suis en forme, alors je vais me balader
\item Si le temps est clément ou que je suis en forme, alors je vais me balader
\item Si le temps est clément, alors je vais me balader et manger une glace
\item Si le temps est clément, alors je vais me balader ou manger une glace
\item Si le temps est clément et que je suis en forme, alors je vais me balader et manger une glace
\item Si le temps est clément et que je suis en forme, alors je vais me balader ou manger une glace
\item Si le temps est clément ou que je suis en forme, alors je vais me balader et manger une glace
\item Si le temps est clément ou que je suis en forme, alors je vais me balader ou manger une glace
\end{enumerate}

\section{Manipulation de formules}

Montrez les équivalences suivantes en utilisant les équivalences remarquables :

\begin{enumerate}
\item $\neg(a \Rightarrow b) \equiv a \wedge \neg b$
\item $a \wedge b \Rightarrow a \equiv \top$
\item $a \Rightarrow a \vee b \equiv \top$
\item $x \Rightarrow (y \wedge z) \equiv (x \Rightarrow y) \wedge (x \Rightarrow z)$
\item $x \Rightarrow (y \vee z) \equiv (x \Rightarrow y) \vee (x \Rightarrow z)$
\item $(a \vee b) \Rightarrow (c \wedge d) \equiv ((a \Rightarrow c) \wedge (a \Rightarrow d)) \wedge ((b \Rightarrow c) \wedge (b \Rightarrow d))$
\item $(a \vee b) \Rightarrow (x \wedge y) \equiv (a \Rightarrow (x \wedge y)) \wedge (b \Rightarrow (x \wedge y))$
\item $(a \wedge b) \Rightarrow (x \vee y) \equiv (a \Rightarrow x) \vee (b \Rightarrow y)$
\item $(x \Rightarrow y) \Rightarrow x \equiv x \vee y$
\item $\neg(x \Rightarrow y) \equiv x \Leftrightarrow \neg y$
\end{enumerate}

\section{Procédure booléenne}

On fournit le pseudo-code de la procédure \texttt{booleenne1}.

\begin{lstlisting}
Procedure booleenne1()
    x <- Entier aleatoire entre 1 et 10
    y <- Entier aleatoire entre 1 et 10
    Si x > y Alors
        Renvoyer FAUX
    Sinon
        Renvoyer VRAI
    FinSi
FinProcedure
\end{lstlisting}

\begin{enumerate}
\item Dans quel cas cette procédure renvoie-t-elle la valeur FAUX ?
\item Dans quel cas cette procédure renvoie-t-elle la valeur VRAI ?
\item Proposez une procédure \texttt{booleenne2} équivalente.
\end{enumerate}

\section{Quelques boucles}

\begin{enumerate}
\item On fournit le pseudo-code de la procédure \texttt{tirage1}. Que peut-on dire une fois cette exécution terminée ? Utilisez une formalisation logique afin de répondre.

\begin{lstlisting}
Procedure tirage1()
    a <- Valeur du tirage d'une piece a pile ou face
    b <- Valeur du tirage d'une piece a pile ou face
    TantQue (a.estPile()&b.estPile())||(a.estFace()&b.estFace()) Faire
        a <- Valeur du tirage d'une piece a pile ou face
        b <- Valeur du tirage d'une piece a pile ou face
    FinTantQue
FinProcedure
\end{lstlisting}

\item On fournit le pseudo-code de la procédure \texttt{tirage2}. Que peut-on dire une fois cette exécution terminée ? Utilisez une formalisation logique afin de répondre.

\begin{lstlisting}
Procedure tirage2()
    a <- Valeur du tirage d'une piece a pile ou face
    b <- Valeur du tirage d'une piece a pile ou face
    TantQue (a.estPile()&b.estPile())||(a.estFace()&!b.estFace()) Faire
        a <- Valeur du tirage d'une piece a pile ou face
        b <- Valeur du tirage d'une piece a pile ou face
    FinTantQue
FinProcedure
\end{lstlisting}

\item Peut-on simplifier ces programmes ?
\end{enumerate}

\section{Quelques conditionnelles}

On fournit le pseudo-code de la procédure \texttt{conditionnelle1}. On donne les variables logiques et propositions suivantes :
\begin{itemize}
\item $a$ : ``$x \leq y$, connaissant $x$ et $y$.''
\item $b$ : ``$x = y + 1$, connaissant $x$ et $y$.''
\item $t$ : ``On affiche 'Titi' ''
\item $o$ : ``On affiche 'Toto' ''
\item $u$ : ``On affiche 'Tutu' ''
\end{itemize}

\begin{lstlisting}
Procedure conditionnelle1()
    x <- Valeur du tirage d'un de a 6 faces
    y <- Valeur du tirage d'un de a 6 faces
    Si x < y Alors
        Afficher "Titi"
    SinonSi x = y + 1 Alors
        Afficher "Toto"
    Sinon
        Afficher "Tutu"
    FinSi
FinProcedure
\end{lstlisting}

\begin{enumerate}
\item Exprimez en fonction de $a$ et $b$ le cas où on affiche 'Titi', c'est-à-dire lorsque la variable $t$ est vraie. Donnez la proposition correspondant à cette formule logique.

\item Exprimez en fonction de $a$ et $b$ le cas où on affiche 'Toto', c'est-à-dire lorsque la variable $o$ est vraie. Donnez la proposition correspondant à cette formule logique.

\item Exprimez en fonction de $a$ et $b$ le cas où on affiche 'Tutu', c'est-à-dire lorsque la variable $u$ est vraie. Donnez la proposition correspondant à cette formule logique.

\item Complétez les conditions dans la procédure suivante, afin qu'elle soit équivalente à \texttt{conditionnelle1}.

\begin{lstlisting}
Procedure conditionnelle2()
    x <- Valeur du tirage d'un de a 6 faces
    y <- Valeur du tirage d'un de a 6 faces
    Si ............. Alors
        Afficher "Tutu"
    SinonSi ............ Alors
        Afficher "Titi"
    Sinon
        Afficher "Toto"
    FinSi
FinProcedure
\end{lstlisting}
\end{enumerate}

\end{document}
